% status: ZERO
% Chapter 29 of the second edition. Plan: docs/STRUCTURE.md. Rules: AUTHORING.md.
% Measurements: research/measurements/2026-09-24-m1-part5.md (reasoning, kvsizes, depth).
\chapter{Reasoning Models and Test-Time Compute}\label{ch:reasoning-workloads}

\begin{ChapterIntent}
Reasoning models generate thousands of tokens the user never sees before
answering, and test-time compute methods sample many candidates and search
among them. The workload becomes dominated by decode, the cache grows with the
output rather than the prompt, and the useful unit of cost is the answer, not
the token. This chapter measures a reasoning trace, shows how parallel
sampling shares prefixes, and follows the consequences for batching,
speculation, caching, scheduling and pricing.
\end{ChapterIntent}

\OpeningSection{Thousands of tokens before the first word}

Ask DeepSeek-R1-0528-Qwen3-8B --- Qwen3-8B post-trained on reasoning
distilled from DeepSeek-R1 \cite{deepseek2025r1qwen3,deepseek2025r1} --- how
many minutes a train takes to go from 3:15~pm to 5:40~pm. On the M1, through
llama.cpp's server, it thought for 2,047 tokens, reached the request's limit of
2,048, and never answered. Across eight such problems --- a multiplication, a
percentage, the divisors of 36 --- it spent on average \textbf{1,237 hidden
tokens} and 166 visible ones per answer, 144 seconds each at about ten tokens
per second, and gave the right number seven times out of eight, six of them in
the requested format (measured, greedy decoding; record
\texttt{2026-09-24-m1-part5}).

Then the same problems with a thinking budget. At 256 tokens, the server forces
the end-of-thinking tag once the budget is spent: 637 tokens and 62 seconds per
answer on average, and all eight numbers right (seven in the requested format).
At a budget of zero, thinking ends before it begins --- and the model reasoned
in its visible answer instead: 1,148 visible tokens per answer, 109 seconds,
seven right and one cut off at the limit in the middle of a calculation. A
budget moved the reasoning more than it removed it. Eight easy problems are an
anecdote about accuracy, but a measurement of cost: every one of those tokens
was a decode step, a cache entry of 144~KiB, and, on a commercial API, a billed
output token.

\Foundations{The answer is not the whole workload}{
A request can compute many tokens before presenting a short final answer. Search can
also produce multiple candidate continuations and then discard most of them. Those
tokens still consume compute, bandwidth and state. Cost per visible token may
consequently be misleading; cost per successfully solved task is often the relevant
outcome. Longer computation does not guarantee a correct answer.
Chapter~\ref{ch:prefill-decode} supplies the phase accounting, and
Chapter~\ref{ch:benchmarking} supplies the measurement contract. Keep generated work,
retained work and delivered output as separate counters.
}

\NapkinSection{Napkin math: cost per answer}

A request with a prompt of $P$ tokens, $R$ reasoning tokens and $A$ answer
tokens costs one prefill of $P$ tokens and $R + A$ decode steps. On hardware
where decode is bandwidth-bound (\Chref{roofline}), each step of one sequence
reads the weights $W$ and a cache that grows by $k$ bytes per token, so the
decode bytes of the whole answer are about
\begin{equation}
(R + A)\,W + k\,\frac{(P + R + A)^{2} - P^{2}}{2},
\label{eq:answer-bytes}
\end{equation}
the second term being the cache reads of every step, which grow with the
square of the output. The time per answer is roughly $(R+A)$ divided by the
decode rate, and the price is $R + A$ output tokens: providers bill hidden
reasoning tokens as output tokens (OpenAI's reasoning guide; Anthropic's
extended-thinking documentation \cite{openai2026reasoning,anthropic2026thinking}).

Put numbers in, with an illustrative 4,000 reasoning tokens, about three
times the opening's average. Qwen3-8B caches 144~KiB per token
(\Chref{shrinking-kv}); a 100-token question answered after 4,000 reasoning
tokens holds 577~MiB of cache
at the end, where the same question answered directly in 40 tokens would hold
20~MiB (derived). On the M1, where this model decodes at 10.5 tokens per second with
an empty context and 8.7 at 4,096 tokens (measured, record
\texttt{2026-09-24-m1-part5}), the 4,000 hidden tokens alone take about seven
minutes. The prompt's prefill, which dominates most chat requests, is a rounding
error here.

\section{Hidden tokens are real tokens}

To the engine, a reasoning token is an output token like any other: sampled,
appended to the cache, counted against the context and the output limit. The
differences are in the protocol around it. OpenAI's API does not return the
reasoning text, but the tokens ``occupy space in the model's context window
and are billed as output tokens'', and the guide recommends reserving at least
25,000 tokens for reasoning and output when starting out
\cite{openai2026reasoning}. Anthropic's manual extended thinking takes a
\texttt{budget\_tokens} of at least 1,024, less than \texttt{max\_tokens},
because thinking counts toward the output limit, and describes the budget as
``a target rather than a strict cap'' \cite{anthropic2026thinking}.

The protocol also decides what happens to reasoning after the answer. Older
Claude models stripped earlier turns' thinking from the context; Claude
Opus~4.5 and models numbered 4.6 and later keep it and bill it as input
\cite{anthropic2026thinking}. OpenAI returns reasoning in encrypted form, to be
passed back between turns or kept server-side with
\texttt{previous\_response\_id} \cite{openai2026reasoning}. Each choice lands on
the engine: kept reasoning is prefix the next turn can reuse from cache
(\Chref{prefix-caching}); stripped reasoning is cache that must be discarded
and a prefix that no longer matches what was computed.

\section{Decode-dominated workloads}

A chat request with a 2,000-token prompt and a 300-token answer spends most of
its tokens in prefill; a reasoning request with a 100-token question and 4,000
tokens of thought spends 98\% of them in decode (derived). Everything Part~\ref{part:serving}
built around prefill changes weight. Chunked prefill (\Chref{chunked-prefill})
matters less when prompts are short; continuous batching
(\Chref{continuous-batching}) matters more, because long outputs finish at
widely different times; and in a disaggregated deployment
(\Chref{disaggregated-serving}) the decode pool must grow relative to the
prefill pool. The throughput that matters is decode throughput at the batch
the cache allows.

\section{Outputs that grow the cache}

A chat server can size a request's cache from its prompt. A reasoning request
admitted with a 100-token prompt may grow to 30,000 tokens, and nobody knows
which one it will be when it is admitted. That breaks two assumptions of the
schedulers in \Chref{scheduling}. Admission by current usage overcommits: a
batch that fits at admission runs out of blocks later, when every request has
grown, and must preempt. And preemption itself costs more: evicting a request
with 20,000 tokens of reasoning means recomputing them or swapping out
gigabytes. The remedies are an output limit per request, admission that
leaves headroom for growth, and a deliberate choice of whom to preempt
(\Chref{scheduling}); and a cache that holds more tokens per byte
(\Chref{shrinking-kv}, \Chref{kv-quantization}, \Chref{hybrid-models}) matters
most for exactly these workloads.

\section{Parallel sampling and shared prefixes}

Test-time compute also scales across samples. Self-consistency samples several
reasoning paths and takes the most common answer; its authors report gains of
17.9 points on GSM8K, 11.0 on SVAMP and 12.2 on AQuA over greedy
chain-of-thought (their measurement, \cite{wang2023selfconsistency}). For the
engine, $n$ samples of one prompt are $n$ sequences that share their prefix:
paged caches store the prompt's blocks once and share them with copy-on-write
(\Chref{paged-kv}, \cite{kwon2023pagedattention}), and prefix caching does the
same across separate requests (\Chref{prefix-caching}). What cannot be shared
is the reasoning itself --- each sample's thousands of tokens are its own --- so
$n$ samples cost almost exactly $n$ times the decode of one.

\section{Search and verification loops}

Sampling and voting is the simplest search. Snell and colleagues compared it
with searching against process-based verifiers and with revising answers
sequentially, and found that allocating test-time compute per prompt by its
difficulty was more than four times as efficient as best-of-$N$ sampling; in a
FLOPs-matched comparison, a smaller model with such a strategy outperformed one
14 times larger on problems where the small model had some success
\cite{snell2024testtime}. For serving, a search loop is a workflow of many
short generations and scoring passes, scheduled like an agent: prefix reuse
across branches, a verifier model beside the generator, and a latency budget
set by the whole loop rather than by one generation.

\section{Budgets and early stopping}

A reasoning budget bounds $R$. s1 made the mechanism explicit as ``budget
forcing'': end the thinking by force when the budget runs out, or extend it by
appending ``Wait'' when the model tries to stop early; with it, a
32-billion-parameter model fine-tuned on 1,000 examples improved from 50\% to 57\% on
AIME24 by thinking longer \cite{muennighoff2025s1}. Qwen3 introduced a
thinking-budget mechanism for its models \cite{yang2025qwen3}, and the commercial APIs
expose budgets (Anthropic's \texttt{budget\_tokens}) or effort levels
(OpenAI's \texttt{none} through \texttt{max}).

In an engine the budget is a sampler. llama.cpp's server implements it as one
at commit \texttt{d006858}: from the start of the thinking block it counts
tokens, and when the budget is spent it forces the end-of-thinking tag by
setting every other token's logit to $-\infty$, waiting first for a complete
UTF-8 character (\texttt{common/reasoning-budget.cpp}); the budget comes from
\texttt{--reasoning-budget} or a request's \texttt{thinking\_budget\_tokens}.
The opening's measurements show what that does on one model: a budget of 256
cut the tokens per answer by 55\% without losing an answer, and a budget of
zero cut them by only 18\%, because the model, denied its thinking block,
thought in the open. A budget bounds the hidden tokens; only the model decides
how many visible ones it needs, which is why the output limit remains the hard
bound on cost.

\section{Speculation for reasoning}

Reasoning is the workload speculation should love --- long outputs at small
batch, decode-bound --- and the one where cheap drafters struggle. On the M1,
prompt lookup proposed no drafts at all for a reasoning model's thoughts, even
when the task was to quote its prompt, and an unmatched small drafter was
accepted less than half the time and made decoding twice as slow
(\Chref{when-speculation-loses}). The candidates are drafters trained on the
target's own distribution --- EAGLE-style heads and multi-token-prediction
modules (\Chref{modern-speculation}); whether they pay on a given reasoning
model is an acceptance rate to measure, not assume.

The answer returned to the user may be short while the engine has evaluated several long candidates. Follow each branch's state separately, and charge the budget for every evaluated branch, including those discarded. Sharing a prefix does not permit branches to share mutable continuation state.

\EngineFigure{sys-reasoning-branches}{Component layout of an illustrative parallel-candidate workload. A common evaluated prefix may be shared read-only, but each branch needs its own mutable continuation state and random stream. Verification or selection is an additional stage; it is not free model work.}{fig:sys-reasoning-branches}

\LedgerSection

Answering with hidden reasoning, or with many sampled candidates:

\begin{ConstraintLedger}
Correctness & buys & Accuracy on hard problems grows with reasoning length and samples, with diminishing returns per token. \\
Latency & spends & Time to the answer is $(R + A)$ decode steps; thousands of steps where a chat answer takes hundreds. \\
Memory capacity & spends & The cache grows with the output; batch size falls as reasoning lengthens, and admission must reserve for growth. \\
Memory bandwidth & moves & Load shifts from prefill to decode, whose bytes grow with the square of the output (Equation~\ref{eq:answer-bytes}). \\
\end{ConstraintLedger}

\LabSection{measure a reasoning trace}

\LabIntent{In \texttt{code/labs/ch29-reasoning-workloads/}, a harness that
runs a reasoning model through an OpenAI-compatible server on a set of problems
with checkable answers, at several thinking budgets and with $n$ parallel
samples, and records per answer: reasoning and answer tokens, time to first
visible token, total time, cache size at the end and correctness. It plots
accuracy against tokens spent and fits the cost model of
Equation~\ref{eq:answer-bytes}. Break it by setting the budget to zero and
checking which answers survive.}

\HappenedSection{serving reasoning models}

In March 2025 DeepSeek published statistics for 24 hours of its V3 and R1
service: 608
billion input tokens, 56.3\% of them served from its on-disk KV cache, and 168
billion output tokens in 24 hours, at an average output speed of 20 to 22
tokens per second per user (the operator's own report,
\cite{deepseek2025inference}). At that speed a reasoning answer of 4,000
tokens takes three minutes. The providers' interfaces show how the industry
responded: budgets and effort levels that let users trade answer quality for
time and price \cite{openai2026reasoning,anthropic2026thinking}, and changing
rules for whether reasoning stays in the context --- Anthropic's newer models
keep earlier thinking, where older ones stripped it
\cite{anthropic2026thinking} --- each of which decides whether the next turn's
prefix is already in cache.

\FrontierSection{2026-09-24}

Reasoning is moving from a fixed mode to a controlled resource: effort levels
and adaptive thinking that decide per request how much to think, budgets
trained into models (Qwen3), and search strategies chosen by estimated
difficulty \cite{snell2024testtime}. For engines, the consequences compound:
decode-dominated batches, caches that grow during generation, and prefix
reuse across turns and branches. The architectures of \Chref{hybrid-models},
whose state does not grow with the output, are attractive for exactly this
reason.

\ExercisesSection

\begin{enumerate}
\item \emph{Derivation.} Derive Equation~\ref{eq:answer-bytes} and find the
  output length at which the cache term exceeds the weights term for Qwen3-8B
  at 4.5 bits per weight.
\item \emph{Prediction.} A server holds 40~GB of free cache memory for a model
  with 144~KiB per token. How many concurrent reasoning requests of 16,000
  tokens fit, and how many chat requests of 2,000?
\item \emph{Implementation.} Add parallel sampling with majority voting to the
  lab, and measure accuracy and cost per answer for $n = 1, 3, 5$.
\item \emph{Falsification.} Find a problem on which a zero thinking budget
  gives the correct answer and a small positive budget gives a wrong one.
\end{enumerate}
